\section{Carta Gantt con línea base y desviaciones}
\label{anexo:gantt}

La Tabla~\ref{tab:gantt-fama} presenta la planificación y la ejecución
de las actividades de F.A.M.A. entre mayo y noviembre de 2026. La línea
base se reconstruyó como propuesta a partir del plan de trabajo de la
memoria. La ejecución se ubicó mediante el anteproyecto, las iteraciones
descritas en el capítulo~7, los registros técnicos y las fechas
confirmadas por el equipo. El inicio del levantamiento, del diseño y
de la redacción se presenta con precisión mensual estimada.

\begin{table}[!htbp]
    \centering
    \caption{Línea base propuesta y ejecución mensual de F.A.M.A.}
    \label{tab:gantt-fama}
    \small
    \setlength{\tabcolsep}{3pt}
    \renewcommand{\arraystretch}{1.13}
    \newcommand{\ganttPlanFama}{\rule{\linewidth}{0.45ex}}
    \newcommand{\ganttRealFama}{\rule{\linewidth}{1.05ex}}
    \begin{tabular}{@{}p{0.35\linewidth}p{0.05\linewidth}*{7}{@{}p{0.08\linewidth}}@{}}
        \hline
        \textbf{Actividad} & &
        \centering\textbf{May} & \centering\textbf{Jun} &
        \centering\textbf{Jul} & \centering\textbf{Ago} &
        \centering\textbf{Sep} & \centering\textbf{Oct} &
        \centering\textbf{Nov}\tabularnewline
        \hline
        Levantamiento de información y requerimientos & P & \ganttPlanFama &  &  &  &  &  &  \\
        & E & \ganttRealFama &  &  &  &  &  &  \\
        \hline
        Diseño de arquitectura, base de datos e interfaz & P &  & \ganttPlanFama & \ganttPlanFama &  &  &  &  \\
        & E &  & \ganttRealFama & \ganttRealFama &  &  &  &  \\
        \hline
        Procesamiento acústico y entrenamiento & P &  &  &  & \ganttPlanFama & \ganttPlanFama &  &  \\
        & E &  &  &  &  & \ganttRealFama &  &  \\
        \hline
        Integración web, servicios en la nube y retroalimentación & P &  &  &  &  & \ganttPlanFama &  &  \\
        & E &  &  &  &  & \ganttRealFama & \ganttRealFama &  \\
        \hline
        Pruebas, ajustes y corrección de errores & P &  &  &  &  & \ganttPlanFama & \ganttPlanFama &  \\
        & E &  &  &  &  & \ganttRealFama & \ganttRealFama & \ganttRealFama \\
        \hline
        Implantación y puesta en marcha local & P &  &  &  &  &  & \ganttPlanFama &  \\
        & E &  &  &  &  &  & \ganttRealFama &  \\
        \hline
        Redacción y corrección de la tesis & P & \ganttPlanFama & \ganttPlanFama & \ganttPlanFama & \ganttPlanFama & \ganttPlanFama & \ganttPlanFama & \ganttPlanFama \\
        & E & \ganttRealFama & \ganttRealFama & \ganttRealFama & \ganttRealFama & \ganttRealFama & \ganttRealFama & \ganttRealFama \\
        \hline
    \end{tabular}
    \par\smallskip
    \parbox{0.97\linewidth}{\footnotesize
    P corresponde a la línea base propuesta (barra fina). E corresponde
    a la ejecución reconstruida (barra gruesa). El último mes de cada
    barra se interpreta al nivel de mes. En redacción, la barra E llega
    al corte del 9 de noviembre de 2026 y la actividad permanece en curso.}
\end{table}

\subsection{Hitos y desviaciones observadas}

El anteproyecto de mayo documentó los perfiles de usuario y los
requerimientos iniciales. La definición oficial de la arquitectura, la
base de datos y la interfaz se situó en julio, según la fecha confirmada
por el equipo. En septiembre se registró la implementación de esa
arquitectura en el código del proyecto.

El 3 de septiembre se registró la primera versión del prototipo
bioacústico en Git. La etapa de desarrollo del procesamiento acústico
y entrenamiento se cerró en septiembre. El calendario propuesto
consideraba su inicio en agosto, por lo que la ejecución se concentró
en septiembre. Los ajustes posteriores de preparación de datos,
evaluación y compatibilidad se reflejan en las actividades de
integración y pruebas.

La integración del backend con el almacenamiento en la nube y la
base de datos tiene un primer registro el 4 de septiembre. La
retroalimentación se incorporó durante ese mes y el corte de integración
se realizó el 24 de octubre. En la comparación mensual, esta actividad
se extendió hasta octubre. La preparación versionada de los datos y
la evaluación con particiones separadas ampliaron las comprobaciones
necesarias para verificar el flujo integrado.

La preparación del entorno oficial comenzó el 23 de octubre y la
actualización del nodo local se verificó el día 25. La puesta en marcha
local quedó situada en octubre. Las pruebas continuaron después de ese
hito y el 9 de noviembre se registró una comprobación adicional de
arranque en Docker con CPU en un entorno aislado. La transferencia
al servidor universitario queda pendiente de programación.

% Fuentes de las fechas y del alcance.
% Anteproyecto de título, versión Avance1, mayo de 2026.
% Memoria actualizada, capítulos 7 a 9.
% Julio y cierre del desarrollo en septiembre, confirmados por el autor.
% Git e888d89 (3 de septiembre), f31c1dd (4 de septiembre).
% Git 0bda3bb (retroalimentación, 27 de septiembre).
% docs/manuales/integracion_dev_20261024.md
% docs/manuales/preparacion_entorno_oficial.md
% docs/manuales/actualizacion_oficial_0004_20261025.md
% docs/receipts/comprobacion_docker_local_20261109.json
